\documentclass[11pt,a4paper]{article}

\usepackage[margin=1in]{geometry}
\usepackage{amsmath}
\usepackage{amssymb}
\usepackage{booktabs}
\usepackage{array}

\title{\vspace{-1.5em}Model Meaning of \texttt{arcDefiningBodies}}
\author{}
\date{}

\begin{document}
\maketitle
\vspace{-2.5em}

\section*{Purpose}

In a multi-arc estimation model, each estimated body's trajectory may be split into
local arcs. The argument \texttt{arcDefiningBodies} specifies which body's local
arc should define the state transition and sensitivity matrix used at a given
epoch.

\section*{Multi-arc state model}

Let body \(b\) have local arcs
\[
    I_{b,k} = [t^s_{b,k}, t^e_{b,k}],
\]
with a separate estimated initial state for each arc,
\[
    x_{b,k,0} = x_b(t^s_{b,k}).
\]
Inside that arc, the state is a function of the local initial state and global
parameters \(p\):
\[
    x_b(t) = \varphi_{b,k}(t; x_{b,k,0}, p),
    \qquad t \in I_{b,k}.
\]

The variational matrix for an active arc is, schematically,
\[
    M_k(t)
    =
    \begin{bmatrix}
        \Phi_k(t) & S_k(t)
    \end{bmatrix}
    =
    \begin{bmatrix}
        \dfrac{\partial x_k(t)}{\partial x_{k,0}} &
        \dfrac{\partial x_k(t)}{\partial p}
    \end{bmatrix}.
\]

For an observation model
\[
    h(t) = h\big(x_{b_1}(t_1), x_{b_2}(t_2), \ldots, p\big),
\]
the partial derivative with respect to estimated arc parameters uses the
variational matrix of the arc that is considered active for the body involved:
\[
    \frac{\partial h}{\partial q}
    =
    \frac{\partial h}{\partial x_b(t)}
    \,
    M_{\alpha(t)}(t),
\]
where \(\alpha(t)\) is the selected arc index. The role of
\texttt{arcDefiningBodies} is to define \(\alpha(t)\).

\section*{Arc selection rule}

Let \(D\) be the set passed through \texttt{arcDefiningBodies}.

\[
\alpha_D(t) =
\begin{cases}
\text{default time-based arc}, & D = \emptyset, \\
\text{arc containing } t \text{ for the body or bodies in } D, & D \neq \emptyset.
\end{cases}
\]

If multiple bodies are provided, they must identify the same active arc. If one
body says ``arc 0'' and another says ``arc 1'', the model request is ambiguous:
there is no single local variational matrix that represents both choices.

\section*{Example with distinct behavior}

Consider two estimated bodies:

\[
\begin{array}{ccl}
\text{Probe } P:  & I_{P,0} = [0,100], & \text{arc index } 0, \\
\text{Relay } R:  & I_{R,1} = [20,30], & \text{arc index } 1.
\end{array}
\]

The relay arc is completely inside the probe arc. At \(t=25\), both local models
exist, but they correspond to different local initial states:

\[
    x_P(25) = \varphi_{P,0}(25; x_{P,0,0}, p),
    \qquad
    x_R(25) = \varphi_{R,1}(25; x_{R,1,0}, p).
\]

The selected variational matrix depends on \(D\):

\begin{center}
\begin{tabular}{>{\ttfamily}lcc}
\toprule
\normalfont Input \(D\) & selected arc & returned matrix \\
\midrule
\{\}        & 1 & \(M_1(25)\) \\
\{P\}       & 0 & \(M_0(25)\) \\
\{R\}       & 1 & \(M_1(25)\) \\
\{P,R\}     & conflict & no unique matrix \\
\bottomrule
\end{tabular}
\end{center}

Thus the same epoch can produce different matrices:
\[
    \texttt{arcDefiningBodies}=\{P\}
    \quad \Rightarrow \quad
    \frac{\partial h}{\partial x_{P,0,0}}
    =
    \frac{\partial h}{\partial x_P(25)} M_0(25),
\]
while
\[
    \texttt{arcDefiningBodies}=\{R\}
    \quad \Rightarrow \quad
    \frac{\partial h}{\partial x_{R,1,0}}
    =
    \frac{\partial h}{\partial x_R(25)} M_1(25).
\]

At \(t=50\), only the probe arc is physically active:
\[
    50 \in I_{P,0},
    \qquad
    50 \notin I_{R,1}.
\]
Passing \(D=\{P\}\) selects \(M_0(50)\). Passing \(D=\{R\}\) identifies no active
relay arc. This is exactly the model-level reason for the input: an observation
partial should use the local arc associated with the body whose state appears in
that partial.

\section*{Interpretation}

\texttt{arcDefiningBodies} is not a dynamical parameter. It is a rule for
choosing the local variational solution in a multi-arc model. It matters when
different estimated bodies have different arc intervals, overlapping arcs, or
observations whose partials are associated with specific link-end bodies.

\end{document}
